\input{_preamble}
\SetLeipzig{obv}{obviative}
% A key that is not ASCII.  Its normal form is a STRING, which under
% pdflatex is the two UTF-8 bytes of the e acute -- and those used to be
% what got typeset, one T1 slot each, so the list printed a key the
% source does not contain.  Spelt three ways on purpose: DECLARED with an
% accent command here, USED as the character below, and exempted from the
% CUSTOMLIST as an accent command again.  All three are one key, or this
% list comes out with a missing entry, an unexplained one, or two.
\SetLeipzig{f\'em}{f\'eminin}
\begin{document}
% The list stands BEFORE every use it reports on: it must be complete all
% the same (that is what the .aux round trip is for).
FRONTLIST
\lpzglist

\ex. \gll Der Hund bellte.\\
     the dog bark.\lpzg{3sg.pst}\\
\glt `The dog barked.'

\ex. \gll Die Frauen schlafen.\\
     \lpzg{def.nom.pl} woman.\lpzg{pl} sleep.\lpzg{prs}\\
\glt `The women sleep.'

% \lpzg inside an \altg stack prints plain, but still counts as used.
\exg. Ich sah \altg{den Mann}{die Frau} gestern.\\
      I see.\lpzg{pst} \altg{the.\lpzg{acc} man}{the.\lpzg{acc} woman} yesterday\\

% declared with \SetLeipzig, so it is explained; NOWHERE is unknown and
% must be dropped from the list without a trace.
Prose: \lpzg{obv} and \lpzg{fém} and \lpzg{NOWHERE}.  \lpzgadd{voc}

INLINELIST \lpzglist[style=inline,sort=false,title={},ignore={3,def,nom}]

CUSTOMLIST
\lpzglist[format={\item[\textsc{#1}] = #2},
          ignore={3,acc,def,f\'em,nom,obv,pl,prs,pst,voc}]
\end{document}
